\subsection{对合代数}

定义 1. --- 设 A 是 C 上的一个代数。A 上的对合是从 A 到 A 的一个映射 \(x \mapsto x^{*}\)，满足：

\[
\begin{array}{c} (x ^ {*}) ^ {*} = x, \quad (x + y) ^ {*} = x ^ {*} + y ^ {*}, \quad (\lambda x) ^ {*} = \overline {{\lambda}} x ^ {*} \\ (x y) ^ {*} = y ^ {*} x ^ {*} \end{array}
\]

对一切 \(x, y \in A\) 及 \(\lambda \in C\) 成立。C 上装备了一个对合的代数称为对合代数。

A 上的对合特别地是从环 A 到反环 A° 上的一个同构。

设 A 是一个对合代数。称 \(x^{*}\) 为 x 的伴随元。A 的在对合下稳定的子集称为自伴的。若 A 有单位元 e，则有 \(e^{*} = e\)；称 (A, e) 为幺对合代数。幺对合代数的元素 u 称为酉元，如果 \(uu^{*} = u^{*}u = e\)，换言之，如果 u 可逆且其逆为 \(u^{*}\)。

元素 \(x \in A\) 称为 Hermite 元，如果 \(x = x^{*}\)；称为正规元，如果 \(xx^{*} = x^{*}x\)。这一术语推广了 A, IX, § 7, 第 3 小节中的术语。每个 Hermite 元都是正规元，每个酉元也都是正规元。A 的 Hermite 元组成的集合 \(A_{h}\) 是 A 的一个实向量子空间。若 x 与 y 是 Hermite 的且可交换，则有 \((xy)^{*} = y^{*}x^{*} = yx = xy\)，故 xy 是 Hermite 的。对每个 \(x \in A\)，A 中的元素 \(xx^{*}\) 与 \(x^{*}x\) 都是 Hermite 的。

若 A = C 装备对合 \(z \mapsto \overline{z}\)，则有 \(A_{h} = R\)。

引理 2. --- 设 A 是一个对合代数，\(x \in \mathrm{A}\)。元素

\[
x _ {1} = \frac {1}{2} (x + x ^ {*}), \qquad x _ {2} = \frac {1}{2 i} (x - x ^ {*})
\]

是 Hermite 的且满足 \(x = x_{1} + ix_{2}\)。若 \(x = y_{1} + iy_{2}\)，其中 \(y_{1}\) 与 \(y_{2}\) 是 Hermite 的，则 \(x_{1} = y_{1}\) 且 \(x_{2} = y_{2}\)。此外，元素 x 是正规元当且仅当 \(x_{1}\) 与 \(x_{2}\) 可交换。

前两个断言由 I, p.~95 的引理 1 得出。经计算有 \(xx^{*} - x^{*}x = 2i(x_{2}x_{1} - x_{1}x_{2})\)，故 x 是正规元当且仅当 \(x_{1}\) 与 \(x_{2}\) 可交换。

设 A 是一个幺对合代数。为使 \(x \in A\) 可逆，必须且只需 \(x^{*}\) 可逆，且此时有 \((x^{*})^{-1} = (x^{-1})^{*}\)。由于 \((x - \lambda e)^{*} = x^{*} - \overline{\lambda}e\) 对一切 \(\lambda \in C\) 成立，我们推出 \(\mathrm{Sp}_{\mathrm{A}}(x^{*}) = \overline{\mathrm{Sp}_{\mathrm{A}}(x)}\)。

设 A 是一个对合代数，\(\widetilde{\mathrm{A}}\) 是由 A 添加单位元而得的代数。\(\widetilde{\mathrm{A}}\) 上存在唯一的拓展 A 的对合的对合，它由 \((\lambda, x)^{*} = (\overline{\lambda}, x^{*})\) 给出，其中 \(\lambda \in \mathbf{C}\)，\(x \in \mathrm{A}\)。若 \(x \in \mathrm{A}\)，则有 \(\mathrm{Sp}_{\mathrm{A}}'(x^{*}) = \overline{\mathrm{Sp}_{\mathrm{A}}'(x)}\)。

设 A 与 B 是对合代数。从 A 到 B 的态射是从 A 到 B 的一个代数态射 \(\varphi\)，使得 \(\varphi(x^{*}) = \varphi(x)^{*}\) 对 A 中一切 \(x\) 成立。A 的恒同映射是对合代数态射。若 C 是对合代数且 \(\pi: B \to C\) 是对合代数态射，则 \(\pi \circ \varphi\) 是对合代数态射。若 \(\varphi\) 是代数同构，则 \(\varphi^{-1}\) 是对合代数态射，此时称 \(\varphi\) 为对合代数同构。

依 I, p.~95 的命题 1，若 A 与 B 是对合代数，则从 A 到 B 的代数态射 \(\varphi\) 是对合代数态射当且仅当 \(\varphi(A_h) \subset B_h\)。我们把 A 的对合子代数称为自伴子代数。A 的中心是一个对合子代数。若 \(A_1\) 是 A 的自伴双边理想，则 A 的对合通过过渡到商在代数 \(A/A_1\) 上定义一个对合，且从 A 到 \(A/A_1\) 上的典范映射是对合代数态射。

设 A 是一个对合代数。A 的根等于反代数的根 (A, VIII, p.~431, 命题 7)，因而是自伴的。

设 A 是一个对合代数。若 M ⊂ A 自伴，则其交换子 M' 是 A 的一个对合子代数。若 \(x \in A\)，则 \(\{x, x^{*}\}\) 的双交换子是包含 x 与 \(x^{*}\) 的对合子代数，且此子代数是交换的当且仅当 x 是正规元 (I, p. 5, 第 5 小节)。

注. — 设 A 是一个对合代数，B 是 A 的一个极大交换对合子代数。那么 B 是一个极大交换子代数。特别地，若 A 是幺的，则代数 B 是全的。

事实上，设 \(x \in A\) 是与 B 可交换的元素。那么 \(x^{*}\) 与 B 可交换。将 x 写成 \(x = x_{1} + ix_{2}\)，其中 \(x_{1}\) 与 \(x_{2}\) 是 Hermite 的；元素 \(x_{1}\) 与 \(x_{2}\) 与 B 可交换（引理 2）。于是 A 的由 B 与 \(x_{1}\) 生成的子代数是交换的且是对合的。从而它等于 B，故 \(x_{1} \in B\)。类似地有 \(x_{2} \in B\)，最终有 \(x \in B\)。

设 A 是一个对合代数。若 \(f\) 是 A 上的一个线性型，则 A 上的映射 \(x \mapsto \overline{f(x^*)}\) 是 A 上的线性型，记作 \(f^*\)。映射 \(f \mapsto f^*\) 是 A' 上的一个半线性对合。我们称 \(f\) 是 Hermite 的，如果 \(f = f^*\)。依 I, p.~95 的引理 1，A 上每个线性型 \(f\) 都有唯一的表示 \(f = f_1 + if_2\)，其中 \(f_1\) 与 \(f_2\) 是 Hermite 的，即 \(f_1 = \frac{1}{2}(f + f^*)\) 与 \(f_2 = \frac{1}{2i}(f - f^*)\)。

为使线性型 \(f\) 是 Hermite 的，必须且只需 \(f\) 在 \(A_h\) 上的限制取实值 (I, p.~95, 命题 1)。映射 \(f \mapsto f|A_h\) 是从 Hermite 线性型组成的实向量空间到实向量空间 \(A_h\) 的对偶向量空间上的一个同构。

特别地，我们用 \(\mathsf{X}^{\prime}(\mathsf{A})_{h}\)（相应地，\(\mathsf{X}(\mathsf{A})_{h}\)）表示 A 的 Hermite 特征标组成的集合（相应地，A 的非零 Hermite 特征标组成的集合）。于是，一个特征标是 Hermite 的，如果它在 \(A_{h}\) 上的限制取实值。

若 A 是交换的且 \(\chi\) 是 A 的特征标，则 \(\chi^{*}\) 是 A 的特征标，且映射 \(\chi \mapsto \chi^{*}\) 是从 \(\mathsf{X}'(\mathsf{A})\) 到 \(\mathsf{X}'(\mathsf{A})\) 上的一个同胚。

引理 3. --- 设 A 是一个交换对合代数。A 到 \(\mathcal{C}_{0}(\mathsf{X}(\mathsf{A}))\) 的 Gelfand 变换是对合代数态射当且仅当 A 的每个特征标都是 Hermite 的。

事实上，说 \(G_{A}\) 是对合代数态射，等于说对一切 \(x \in A\) 与 \(\chi \in \mathsf{X}(A)\) 有

\[
\chi (x ^ {*}) = \mathcal {G} _ {\mathrm{A}} (x ^ {*}) (\chi) = \overline {{\mathcal {G} _ {\mathrm{A}} (x) (\chi)}} = \overline {{\chi (x)}},
\]

也就是说，每个 \(\chi\) 都是 Hermite 的。

例. --- 1) 设 A 是集合 X 上复值函数组成的代数。映射 \(f \mapsto \overline{f}\) 是 A 中的一个对合。X 上有界函数组成的子代数是 A 的对合子代数。若 X 是局部紧拓扑空间，则子代数 \(\mathcal{C}(X), \mathcal{C}_b(X), \mathcal{C}_0(X)\) 与 \(\mathcal{K}(X)\) 都是 A 的对合子代数。

\begin{enumerate}
\def\labelenumi{\arabic{enumi})}
\setcounter{enumi}{1}
\item
  设 X 是局部紧拓扑空间，\(\mu\) 是 X 上的一个正测度。映射 \(f \mapsto \overline{f}\) 是代数 \(\mathcal{L}^{\infty}(X, \mu)\) 上的一个对合；它通过过渡到商在幺代数 \(\mathrm{L}^{\infty}(X, \mu)\) 上诱导出一个对合。
\item
  设 E 是一个复 Hilbert 空间。在 Banach 代数 \(\mathcal{L}(\mathrm{E})\) 上，映射 \(x \mapsto x^{*}\) (EVT, V, p.~37, 命题 1) 是一个对合。
\item
  设 G 是一个局部紧群。设 \(\mathcal{M}^{1}(G)\) 是 G 上有界复测度组成的 Banach 代数 (I, p.~19, 例 6)。
\end{enumerate}

从 G 到 G 的映射 \(x \mapsto x^{-1}\) 把每个测度 \(\mu \in \mathcal{M}^{1}(G)\) 映为测度 \(\check{\mu} \in \mathcal{M}^{1}(G)\) (INT, VII, p.~12, 公式 (13))。我们用 \(\mu^{*}\) 表示 \(\check{\mu}\) 的复共轭测度。映射 \(\mu \mapsto \check{\mu}\) 是从 Banach 代数 \(\mathcal{M}^{1}(G)\) 到 Banach 代数 \(\mathcal{M}^{1}(G^{\circ})\) 上的一个等距同构 (INT, VIII, §3, 第 1 小节, 命题 7 的推论)；于是 \(\mu \mapsto \mu^{*}\) 是 Banach 代数 \(\mathcal{M}^{1}(G)\) 的一个等距对合。

关于一个 Haar 测度具有密度的有界测度组成的集合 A 是 \(\mathcal{M}^{1}(\mathrm{G})\) 的一个在对合下稳定的闭子代数（参见 INT, VIII, §4, 第 5 小节）；它与 Haar 测度的选取无关。

设 \(\nu\) 是 G 上的一个左 Haar 测度，用 \(\Delta\) 表示 G 的模函数。我们给 \(\mathrm{L}^1 (\mathrm{G},\nu)\) 装备乘积 \((f,g)\mapsto f*\nu g\) 与对合 \(f\mapsto f^{*} = \widetilde{f}\cdot \Delta^{-1}\)，其中 \(\widetilde{f} (x) = \overline{f(x^{-1})}\) 对一切 \(x\in \mathbf{G}\) 成立。于是映射 \(f\mapsto f\cdot \nu\) 是从对合代数 \(\mathrm{L}^1 (\mathrm{G},\nu)\) 到 A 上的一个同构。这个同构是等距的。特别地，\(\mathrm{L}^1 (\mathrm{G},\nu)\) 恒等同于 \(\mathcal{M}^1 (\mathrm{G})\) 的一个对合子代数。

\begin{enumerate}
\def\labelenumi{\arabic{enumi})}
\setcounter{enumi}{4}
\tightlist
\item
  设 U 是 C 的一个在复共轭下稳定的开子集。考虑 U 上复值全纯函数组成的代数 \(\mathcal{O}(\mathrm{U})\)。对每个函数 \(f \in \mathcal{O}(\mathrm{U})\)，映射 \(f^{*}: z \mapsto \overline{f(\overline{z})}\) 是 U 上的一个全纯函数。映射 \(f \mapsto f^{*}\) 是 \(\mathcal{O}(\mathrm{U})\) 上的一个对合。
\end{enumerate}

类似地，设 S 是 C 的一个在复共轭下稳定的紧子集。考虑在 S 的邻域中的复值全纯函数芽组成的代数 \(\mathcal{O}(S)\)。映射 \(f \mapsto f^{*}\) 是 \(\mathcal{O}(S)\) 上的一个对合。

在 I, p.~85 第 13 小节中定义的子代数 \(\mathcal{O}_{\mathbf{R}}(\mathrm{U})\)（相应地，子代数 \(\mathcal{O}_{\mathbf{R}}(\mathrm{S})\)）是 \(\mathcal{O}(\mathrm{U})\)（相应地，\(\mathcal{O}(\mathrm{S})\)）的 Hermite 元组成的集合。

